\centerline{\Huge Памятка 1.0}
\centerline{\Huge }

\textbf{Основное тригонометрические тождество.} $\cos^2\alpha + \sin^2\alpha = 1$ при любом действительном $\alpha$

\textbf{Чётность косинуса.} $\cos(-\alpha) = \cos(\alpha)$

\textbf{Нечётность синуса.} $\sin (-\alpha) = - \sin(\alpha) $
 
\textbf{Синус суммы/разности.} $sin(\alpha \pm \beta) = sin \, \alpha \, cos \, \beta  \pm sin \, \beta \, cos \, \alpha $

\textbf{Косинус суммы/разности} $cos(\alpha \pm \beta) = cos \, \alpha \, cos \, \beta  \mp sin \, \alpha \, sin \, \beta $

\textbf{Формулы приведения аргумента.} 

$$
\sin(\frac{\pi}{2}- \alpha) = \cos \alpha \qquad \cos(\frac{\pi}{2} - \alpha) = -\sin\alpha
$$

$$
\sin(\pi - \alpha) = \sin\alpha, \qquad \cos(\pi - \alpha) = - \cos\alpha
$$

\textbf{Формулы двойного аргумента.} $\sin2\alpha = 2\sin\alpha \cos\alpha, \cos2\alpha = 2\cos^2\alpha - 1 = 1 - 2\sin^2\alpha$

\textbf{Теорема синусов.} $\frac{a}{\sin \alpha} = \frac{b}{\sin \beta} = \frac{c}{\sin \gamma} = 2R$

\textbf{Теорема косинусов.} $a^2 = b^2 + c^2 - 2bc  \cos  \alpha$

\textbf{Формула Стюарта.} Пусть $D$ лежит на стороне $BC$ треугольника $ABC$. Тогда верно следующее тождество: $AD^2 \cdot BC = AB^2 \cdot CD + AC^2 \cdot BD - BC \cdot CD \cdot BD.$ 

\textbf{Альтернативная форма записи формулы Стюарта.} Пусть $D$ лежит на стороне $BC$ треугольника $ABC$ и $BD = x,\ CD = y$. Тогда $AD^2 = \frac{c^2y}{a} + \frac{b^2x}{a} - xy$

Формулы площади треугольника: $S = pr = \frac{1}{2}ah_a = \frac{1}{2}bc \, sin \, \alpha = \frac{abc}{4R} = \sqrt{p(p-a)(p-b)(p-c)}$

\textbf{Длина медианы и биссектрисы.} $m_a = \frac{1}{2}\sqrt{2b^2 + 2c^2 - a^2}$ и $l_a = \sqrt{bc - xy}$, где $x, y$ отрезки на которые делит биссектрисы противоположную сторону.

\textbf{Теорема Чевы в синусной форме.} Дан треугольник $ABC$. Прямые $AA_1, BB_1, CC_1$ пересекаются в одной точке тогда и только тогда, когда

$$
	\frac{\sin \angle(AB, AA_1)}{\sin \angle(AA_1, AC)} = 
	\frac{\sin \angle(CA, CC_1)}{\sin \angle(CC_1, CB)} =
	\frac{\sin \angle(BC, BB_1)}{\sin \angle(BB_1, BA)} = 
	1
$$
